% short/02_data.tex: the app and the inference page, the protocol and the app's markers, participants (Figure 3, the
% session photo), what was kept (598 trials), what a recording looks like (Figure 4, raw signals). PD01: diagnosis
% stated, age band only, no sex or gendered pronoun, no height, device or dates. The video footnote is the one of the
% full report's data section, unchanged. The inference page finds each movement without the start beep (lane m4's
% detector, never evaluated with the final model): said in one plain sentence once the markers are defined.
\section{Data}\label{sec:data}

\paragraph{The app.} We built the app that recorded the data, NoseTouch: a web app that runs in the phone's browser
on iPhones (iOS) and Android phones, with nothing to install. It guides the participant hands-free with spoken and
visual prompts, records the accelerometer and the gyroscope, and uploads every recorded movement. We also built a
second page that runs the final model of this report in the phone's browser and shows which hand it predicts for a
new person's movement (the inference page, \predictorlink).

\paragraph{The protocol.} A participant holds the phone still at arm's length, with the screen facing them. At a
start beep they bring the phone to touch the nose, hold it there until a double beep, and stretch the arm out again
(Figure~\ref{fig:gesture}); one such movement is a trial. The app asks for one hand at a time, in short blocks, and
varies the grip (the phone upright, or sideways with its top to the left or to the right), the speed (natural, fast or
slow) and, in some trials, the movement itself (for example a curved path or a pause halfway). While recording, the
app detects when the phone starts to move after the beep and when it comes to rest at the nose (it needs this moment
to play the double beep), and it stores these two moments with the recording; we call them the app's markers. The
inference page has no start beep: it must find each movement in the sensor stream by itself, a step that this report
does not evaluate (every accuracy in it uses the app's markers), so its answers may be less accurate. A participant's
visit, one session with the app, took a median of \R{describe/visit_minutes_median}~minutes.

\begin{figure}[tbp]
\centering
\repoimage[0.64\linewidth]{photo}{figures/photo/data_collection.jpg}
\caption{A data-collection session (cropped so that no face is visible). Each participant follows the prompts on
their own phone, starting from arm's length (the phone on the left asks for the left hand at natural speed); the
sweets on the table were the compensation.}
\label{fig:photo}
\end{figure}

\paragraph{Participants.} The \R{describe/n_persons} participants, P01 to P20, were Bar-Ilan University students,
most of them master's students, each offered sweets for their time (Figure~\ref{fig:photo}; a video shows a
session\footnote{\raggedright A video of a data-collection session:
\url{https://drive.google.com/file/d/1V9PF623HkJ8MjR6loJhUHnDkCx60W-Jl/view?usp=sharing}.}). They sat or stood while recording
(\R{x1/ref_sessions_sitting} and \R{x1/ref_sessions_standing} sessions), and \R{describe/n_persons_multi_session} of
them recorded a second or third session, so there are \R{describe/sessions_included} sessions in all. They were
\R{describe/n_female} women and \R{describe/n_male} men with a mean age of \R{describe/age_mean} (SD
\R{describe/age_sd}) years, from \R{describe/age_min} to \R{describe/age_max}. The mean height was
\R{describe/height_mean_female} (SD \R{describe/height_sd_female})~cm for the women and \R{describe/height_mean_male}
(SD \R{describe/height_sd_male})~cm for the men (\R{describe/height_n_female} of the \R{describe/n_female} women and
\R{describe/height_n_male} of the \R{describe/n_male} men gave their height). They also reported their handedness
(\R{describe/n_right_handed} right-handed, \R{describe/n_left_handed} left-handed) and the hand they usually hold a
phone in; \R{describe/persons_android} recorded on an Android phone and \R{describe/persons_ios} on an iPhone. One
more person, PD01, who has a diagnosis of Parkinson's disease and is in their fifties, recorded with the same app.
PD01 was never used to build or test any model and is described separately in Section~\ref{sec:secondary}.

\paragraph{What we kept.} Of the \R{describe/attempts_total} tries recorded (including those the app rejected), we
kept \R{describe/primary_trials} trials: those that the app accepted (among other checks, the phone
started to move after the start beep and came to rest at the nose within the time limits, and the participant did
not ask to redo the trial), that lasted at least \R{describe/min_duration_s}~s, in which the phone clearly turned (a
peak rotation speed of at least \R{describe/min_peak_gyro_dps} degrees per second), and whose stream of samples had no
gap longer than \R{describe/max_gap_ms}~ms around the movement (phones sometimes stop sending samples for a moment).
We did not use a warm-up trial at the start of each visit, in which the participant chose the hand, nor a few
discarded sessions: one that a participant stopped and later repeated, and the failed recordings of one more
volunteer. The label of a trial is the hand the app asked for. The \R{describe/primary_trials} trials are
\R{describe/primary_left} left-hand and \R{describe/primary_right} right-hand trials, a median of
\R{describe/primary_per_person_median} per participant (from \R{describe/primary_per_person_min} to
\R{describe/primary_per_person_max}).

\begin{figure}[tbp]
\centering
\begin{minipage}[c]{0.49\linewidth}\centering\genfig{core/S2_22_raw_example}\end{minipage}\hfill
\begin{minipage}[c]{0.48\linewidth}
\caption{Two trials of P15, one with each hand (the same Android phone, upright grip, natural speed), exactly as
the phone recorded them, about \R{describe/rate_hz_median_whole} samples per second with no smoothing (here the
colours mark the axes, not the hands). Top: the gyroscope,
the rotation speed about each of the phone's axes ($x$ across the screen, $y$ along it, $z$ out of it); bottom: the
accelerometer along the same axes, gravity included (in this grip the $y$ axis points up, so it reads about
9.8~m/s$^2$ even at rest). Grey band: the movement to the nose,
between the app's markers; beige band: the arm going back. Time runs from the start beep.}
\label{fig:raw}
\end{minipage}
\end{figure}

\paragraph{What a recording looks like.} Both sensors measure along the phone's own three axes: $x$ across the screen
(towards its right edge), $y$ along the screen (towards its top) and $z$ out of the screen. The gyroscope gives how
fast the phone rotates about each axis, in degrees per second; the accelerometer gives the acceleration along each
axis, in m/s$^2$, and also feels gravity, so an Android phone held still and upright reads about 9.8~m/s$^2$ on its
$y$ axis (Section~\ref{sec:processing} explains iPhones). Figure~\ref{fig:raw} shows one left-hand and one right-hand
trial of the same participant exactly as recorded: six noisy channels per trial (three rotation speeds and three
accelerations), in which the two trials differ, but nothing tells the eye which hand moved the phone.
